\section{Data Recipe：数据不是库存，而是闭环}

有了整机和供应链，下一步才谈得上数据。这里的顺序不能倒过来：机器人数据不是从网上下载的文本，而是在硬件、场景、任务和客户反馈中被生产出来的。

EP132 与 EP134 的“数据综述”可以互相照亮。高继扬这里讲的 Data Recipe，核心不是把数据当产品卖，而是把整机、场景、任务、采集、训练、评测和开源生态连成可迭代飞轮。

\begin{figure}[H]
\centering
\includegraphics[width=0.96\textwidth]{robot-data-recipe-flywheel.png}
\caption{机器人 Data Recipe 飞轮：整机部署、数据采集、模型训练、现场价值、错误回流和开放共享。}
\end{figure}

\begin{knowledgebox}{读图：为什么 Data Recipe 不是卖数据}
图中的数据来自真实硬件和场景。整机进入开发者或生产力场景后产生轨迹、失败和反馈；这些数据训练 VLA、基础模型和后训练工具；模型创造客户价值后带来更多场景与反馈。数据不是静态库存，而是由产品和客户持续生成的学习燃料。
\end{knowledgebox}

\subsection{万台出货为什么重要}

高继扬把“在生产力场景做出万台出货量”称为当下的重要 bet。这个目标不是单纯销量数字，而是数据、供应链、质量和客户价值的综合验证。万台意味着硬件能交付，客户愿意用，维护体系能跟上，数据回流能规模化。

\begin{knowledgebox}{术语消化：机器人数据闭环}
\begin{tabularx}{\textwidth}{p{0.23\textwidth}p{0.32\textwidth}X}
\toprule
术语 & 解决的问题 & 本期中的含义 \\
\midrule
Data Recipe & 数据来源、筛选、训练和反馈配方 & 决定真实数据如何变成模型能力。 \\
Demo in Video & 视频里看起来成立的演示 & 容易塑造过高预期。 \\
Demo in Office & 在公司办公室现场展示 & 比视频更强，但仍是可控环境。 \\
Demo in the Wild & 在客户、展会、跨国家场景中部署演示 & 更接近真实泛化能力。 \\
生产力场景 & 客户用机器人完成实际工作 & 数据和商业价值闭环的来源。 \\
\bottomrule
\end{tabularx}
\end{knowledgebox}

\subsection{开源与分享}

星海图开源数据集和模型，并向客户分享数据。高继扬强调数据不是主要业务，但乐于分享。这可以理解为生态策略：让更多开发者和客户验证任务、暴露需求、形成反馈，从而扩大基础模型和工具链的使用场景。

\subsection{本章小结}

Data Recipe 的本质是把整机、客户和模型训练连起来。机器人公司真正稀缺的不是“有一批数据”，而是持续产生高价值数据、评测和反馈的机制。

